\documentclass[a4paper,twoside]{article}

\usepackage{morewrites}

\usepackage[svgnames,dvipsnames,x11names]{xcolor}
\usepackage{graphicx}
\usepackage[english]{babel}
\usepackage[colorlinks=true, linkcolor=NavyBlue, urlcolor=NavyBlue, citecolor=NavyBlue]{hyperref}
\usepackage[a4paper]{geometry}
\usepackage{ts-fonts}
\usepackage{booktabs}
\usepackage{tabularx}
\usepackage{amsmath}
\usepackage{float}
\usepackage{seqsplit}
\newcolumntype{Y}{>{\centering\arraybackslash}X}
\newcolumntype{Z}{>{\raggedleft\arraybackslash}X}
\usepackage{ts-code}

\usepackage{lastpage}
\usepackage{refcount}
\makeatletter
\def\ps@plain{
  \let\@oddhead\@empty
  \let\@evenhead\@empty
  \def\@oddfoot{
    \hfil
    \ifnum\getpagerefnumber{LastPage}>1\relax
      \thepage
    \fi
    \hfil}
  \let\@evenfoot\@oddfoot
}
\makeatother

\title{Emoji Support}
\author{}\date{}
\setlength{\emergencystretch}{3em}
\widowpenalties 3 10000 10000 150
\clubpenalties  3 10000 10000 150
\displaywidowpenalty 10000
\renewenvironment{abstract}{
  \small
  \begin{center}{\bfseries\abstractname}\end{center}
  \begin{quotation}
}{\end{quotation}}
\makeatletter
\newcount\ts@affilcnt
\newcount\ts@loopcnt
\renewcommand{\thanks}[1]{
  \stepcounter{footnote}
  \protected@xdef\@thefnmark{\thefootnote}
  \@footnotemark
  \global\advance\ts@affilcnt\@ne
  \expandafter\xdef\csname ts@affilM\the\ts@affilcnt\endcsname{\@thefnmark}
  \expandafter\gdef\csname ts@affilT\the\ts@affilcnt\endcsname{#1}
}
\newcommand\tsreplayfn{
  \ts@loopcnt\z@
  \loop
    \advance\ts@loopcnt\@ne
    \ifnum\ts@loopcnt>\ts@affilcnt\else
      \xdef\@thefnmark{\csname ts@affilM\the\ts@loopcnt\endcsname}
      \@footnotetext{\csname ts@affilT\the\ts@loopcnt\endcsname}
  \repeat
}
\makeatother

\begin{document}

\twocolumn[{
\maketitle
}]
\tsreplayfn

\thispagestyle{plain}
\section{Introduction}\label{introduction}

TeXSmith renders emoji as glyphs when you pick a font flavour:

\begin{code}{yaml}{}{}
\begin{Verbatim}[commandchars=\\\{\},breaklines, breakanywhere, commandchars=\\\{\}]
\PY{n+nt}{press}\PY{p}{:}
\PY{+w}{  }\PY{n+nt}{fonts}\PY{p}{:}
\PY{+w}{    }\PY{n+nt}{emoji}\PY{p}{:}\PY{+w}{ }\PY{l+lScalar+lScalarPlain}{black}
\end{Verbatim}
\end{code}
You can choose among four built-in options:

\begin{description}
\item[{ \texttt{black} }] OpenMoji Black (default).
\item[{ \texttt{color} }] Noto Color Emoji.
\item[{ \texttt{twemoji} }] Use the \texttt{twemoji} package as fallback.
\item[{ \texttt{artifact} }] Download emoji as images using Twemoji.

\end{description}
Any other name is treated as a custom font family to load directly.

Engines:

\begin{itemize}
\item{} LuaLaTeX relies on \texttt{luaotfload} to add the emoji font as a fallback.
\item{} XeLaTeX/Tectonic use \texttt{ucharclasses} to automatically switch to the emoji font on the U+1F000--U+1FAFF range.
You can type emoji directly in Markdown or LaTeX source.

\end{itemize}
\newpage
\section{Examples}\label{examples}

\begin{center}
\begin{tabularx}{\linewidth}{>{\raggedright\arraybackslash}X>{\raggedright\arraybackslash}X}
\toprule
\textbf{Emoji} & \textbf{Description} \\

\midrule
\texsmithEmoji{😊} & Smiling face with smiling eyes \\
\texsmithEmoji{🚀} & Rocket \\
\texsmithEmoji{🍕} & Pizza \\
\texsmithEmoji{🎉} & Party popper \\
\texsmithEmoji{🐍} & Snake \\
\texsmithEmoji{🌍} & Globe showing Europe-Africa \\
\texsmithEmoji{💻} & Laptop computer \\
\texsmithEmoji{📚} & Books \\
\texsmithEmoji{🎨} & Artist palette \\
\texsmithEmoji{👽} & Alien \\
\texsmithEmoji{👋} & Waving hand \\
\texsmithEmoji{🤖} & Robot \\
\texsmithEmoji{🦄} & Unicorn \\
\texsmithEmoji{🧠} & Brain \\
\texsmithEmoji{🛸} & Flying saucer \\
\texsmithEmoji{🛰️} & Satellite \\
\texsmithEmoji{🐙} & Octopus \\
\texsmithEmoji{📝} & Memo \\
\texsmithEmoji{📋} & Note \\
\texsmithEmoji{⭐} & Star \\
\texsmithEmoji{✅} & Check mark \\
\texsmithEmoji{❌} & Cross mark \\
\texsmithEmoji{🧪} & Experiment \\
\texsmithEmoji{💡} & Light Bulb \\

\bottomrule
\end{tabularx}
\end{center}

\end{document}
